% Part: computability
% Chapter: computability-theory
% Section: introduction

\documentclass[../../../include/open-logic-section]{subfiles}

\begin{document}

\olfileid[pt]{cmp}{thy}{int}
\olsection{Introdução}

O ramo da lógica conhecido como \emph{teoria da computabilidade} trata
de questões relacionadas à computabilidade, ou computabilidade relativa,
de funções e conjuntos. O fato de essa área ter sido conhecida como
\emph{teoria da recursão}, sendo hoje comuns as duas denominações,
atesta a influência de Kleene.

Chamemos uma função~$f\colon \Nat \pto \Nat$ de \emph{parcialmente
  computável} se ela puder ser calculada em algum modelo de computação.
Se $f$ for total, diremos simplesmente que $f$~é \emph{computável}.
Uma relação~$R$ cuja função característica~$\Char{R}$ é computável
também é chamada de computável. Se $f$ e~$g$ são funções parciais,
escreveremos $f(x) \fdefined$ para indicar que $f$ está definida
em~$x$, O.s.\@, que $x$~pertence ao domínio de~$f$; e
$f(x) \fundefined$ para indicar o contrário, O.s.\@, que $f$ não está
definida em~$x$. Usaremos $f(x) \simeq g(x)$ para indicar que ou
$f(x)$ e $g(x)$ são ambas indefinidas, ou ambas estão definidas e são
iguais.

É possível explorar a teoria da computabilidade sem fazer referência a
um modelo específico de computação. Para isso, mostra-se que existe uma
função parcialmente computável universal~$\fn{Un}(k, x)$. Isso nos
permite enumerar as funções parcialmente computáveis. Adotaremos a
notação~$\cfind{k}$ para designar a $k$-ésima função parcialmente
computável unária, definida por
$\cfind{k}(x) \simeq \fn{Un}(k, x)$. (Kleene usava $\{ k \}$ para
esse fim, mas essa notação tem sido menos usada recentemente.) De modo
um pouco mais geral, podemos enumerar uniformemente as funções
parcialmente computáveis de aridades arbitrárias, e usaremos
$\cfind{k}[n]$ para designar a $k$-ésima função recursiva parcial
$n$-ária.

Se $f(\vec x, y)$ é uma função total ou parcial, então $\umin{y}{f
(\vec x, y)}$ é a função de~$\vec x$ que devolve o menor~$y$ tal que
$f(\vec x, y) = 0$, supondo que todos os valores $f(\vec x, 0)$,
\dots, $f(\vec x, y-1)$ estejam definidos; se não houver tal $y$,
$\umin{y}{f (\vec x, y)}$ será indefinida. Se $R(\vec x, y)$ é uma
relação, define-se $\umin{y}{R(\vec x, y)}$ como o menor~$y$ tal que
$R(\vec x, y)$ é verdadeira; em outras palavras, o menor~$y$ tal que
$1 \tsub \Char{R}(\vec x, y) = 0$.

Para mostrar que uma função é computável, há
duas maneiras de proceder:
\begin{enumerate}
\item Rigorosamente: descrever explicitamente uma máquina de Turing ou
  uma função recursiva parcial e mostrar que ela calcula a função que
  se tem em mente;
\item Informalmente: descrever um algoritmo que a calcula e apelar à
  tese de Church.
\end{enumerate}
Não há uma linha divisória nítida entre as duas maneiras; uma descrição
detalhada de um algoritmo deve fornecer informação suficiente para que
fique relativamente claro como se poderia, em princípio, projetar a
máquina de Turing ou a sequência de definições recursivas parciais
adequada. Definições plenamente rigorosas dificilmente seriam
esclarecedoras, e tentaremos encontrar um meio-termo apropriado entre
essas duas abordagens; em suma, tentaremos encontrar demonstrações
intuitivas, mas rigorosas, de que seria possível obter as definições
precisas.

\end{document}

